\documentclass[a4paper,12pt]{article}
\usepackage[T2A]{fontenc}
\usepackage[utf8]{inputenc}
\usepackage[russian]{babel}
\usepackage{amsmath,amssymb,mathtools}
\usepackage{geometry}
\geometry{left=24mm,right=24mm,top=23mm,bottom=24mm}
\usepackage{microtype,parskip}
\usepackage{xcolor,tabularx,booktabs,tikz,hyperref}
\usetikzlibrary{calc}
\definecolor{accent}{HTML}{008080}
\definecolor{darkaccent}{HTML}{004D4D}
\hypersetup{colorlinks=true,linkcolor=darkaccent,pdftitle={Ревью домашней работы — Ксения Клыкова}}
\pagestyle{plain}
\begin{document}
\hypersetup{pageanchor=false}
\begin{titlepage}
\centering\vspace*{30mm}
{\large\color{darkaccent} МАТЕМАТИКА\par}
\vspace{17mm}{\LARGE Ревью домашней работы\par}
\vspace{12mm}{\Large\color{accent} Теорема Менелая и сечения пирамиды\par}
\vspace{18mm}{\normalsize Геометрия. Анализ рукописных решений\par}
\vfill
{\large Ксения Клыкова\par}
\vspace{12mm}{\normalsize 21 ноября 2026 г.\par}
\vspace{22mm}{\small Лёвин Артём Александрович\par}
{\small эксклюзивно для Ксении Клыковой\par}
\vspace{16mm}
\begin{tikzpicture}
\draw[accent,line width=.8pt](0,0)..controls (1,.5) and (2.4,-.45)..(3.6,0)..controls (4.9,.5) and (6.2,-.5)..(7.3,0);
\end{tikzpicture}
\end{titlepage}
\hypersetup{pageanchor=true}
\setcounter{page}{1}
\newpage
\begin{center}{\Large\color{darkaccent} Сводная таблица}\end{center}
\vspace{12mm}
\begingroup\small\setlength{\tabcolsep}{3pt}
\begin{tabularx}{\linewidth}{@{}p{10mm}p{27mm}>{\raggedright\arraybackslash}X p{19mm}p{22mm}@{}}
\toprule
\textbf{№}&\textbf{Тема}&\textbf{Суть ошибки}&\shortstack[l]{\textbf{Ваш}\\\textbf{ответ}}&\shortstack[l]{\textbf{Верный}\\\textbf{ответ}}\\
\midrule
\hyperref[task:1]{1}&Менелай&Неверно сокращена дробь $15/48$; дальнейший расчёт искажён.&$12:5$&$32:5$\\
\bottomrule
\end{tabularx}
\endgroup
\par\vspace{9mm}
В таблице приведены только задания, требующие исправления.
Нажмите на номер для перехода к подробному разбору.
\clearpage
\thispagestyle{empty}
\vspace*{.35\textheight}
\begin{center}{\Huge\color{darkaccent} Разбор ошибок}\end{center}
\clearpage
\phantomsection\label{task:1}
\section*{Задание № 1. Теорема Менелая}
\textbf{Текст задания.}
В пирамиде $SABC$ точки $P\in AS$, $Q\in CS$, $N\in CB$ удовлетворяют
\[
AP:PS=6:5,\quad CQ:QS=3:8,\quad CN:NB=2:1.
\]
В каком отношении плоскость $(PQN)$ делит ребро $AB$?

Обозначим $M=AB\cap(PQN)$. Продолжение прямой $PQ$ пересекает
прямую $AC$ в точке $F$. Точки $M,N,F$ коллинеарны,
поскольку прямая $NF$ является линией пересечения плоскостей $(PQN)$ и $(ABC)$.

\textbf{Ваше решение.}
Правильно применена теорема Менелая в треугольнике $SAC$:
\[
\frac{AP}{PS}\cdot\frac{SQ}{QC}\cdot\frac{CF}{FA}=1,\qquad
\frac65\cdot\frac83\cdot\frac{CF}{FA}=1.
\]
Получено правильное промежуточное значение $\frac{CF}{FA}=\frac{15}{48}$.
Затем выполнен переход $\frac{15}{48}=\frac56$.
Используя ошибочную дробь в следующем уравнении Менелая,
Вы получили ответ $AM:MB=12:5$.

\textbf{Ваш ответ:} $12:5$.

\textcolor{red}{\textbf{Ошибка:}} неверное сокращение дроби.

\textbf{Где именно возникает ошибка.}
Последний правильный шаг: $\displaystyle\frac{CF}{FA}=\frac{15}{48}$.
Первый неверный шаг: $\displaystyle\frac{15}{48}=\frac56$.

\textbf{Почему это неверно.}
Сокращение дроби означает деление числителя и знаменателя
на один и тот же общий ненулевой делитель. Общий делитель $15$ и $48$ равен $3$:
\[
\frac{15}{48}=\frac{15:3}{48:3}=\frac5{16}.
\]
Действительно, $48:3=16$, а не $6$. Проверка крест-накрест
также показывает различие дробей: $15\cdot6=90$,
тогда как $48\cdot5=240$.
Ошибочное значение $5/6$ переходит в следующие вычисления и меняет ответ.

\textcolor{blue!60!black}{\textbf{Ключевая идея:}}
сохраните геометрическую конструкцию и оба применения теоремы Менелая,
исправив сокращение и все последующие вычисления.

\textbf{Исправление от последнего правильного шага.}
\[
\frac{CF}{FA}=\frac{15}{48}=\frac5{16}.
\]
В треугольнике $ABC$ для точек $M\in AB$, $N\in BC$ и
$F$ на продолжении $AC$ теорема Менелая даёт
\[
\frac{AM}{MB}\cdot\frac{BN}{NC}\cdot\frac{CF}{FA}=1.
\]
Из $CN:NB=2:1$ следует $BN/NC=1/2$. Поэтому
\[
\frac{AM}{MB}\cdot\frac12\cdot\frac5{16}=1,\quad
\frac{AM}{MB}\cdot\frac5{32}=1,\quad
\boxed{\frac{AM}{MB}=\frac{32}{5}}.
\]

\textbf{Полное правильное решение.}
\begin{enumerate}
\item Проведите $PQ$ до пересечения с $AC$ в точке $F$.
По Менелаю в $SAC$:
\[
\frac{AP}{PS}\cdot\frac{SQ}{QC}\cdot\frac{CF}{FA}=1.
\]
\item Подставьте данные и выразите отношение:
\[
\frac65\cdot\frac83\cdot\frac{CF}{FA}=1
\quad\Longrightarrow\quad
\frac{CF}{FA}=\frac{5\cdot3}{6\cdot8}=\frac5{16}.
\]
\item Пусть $M=AB\cap(PQN)$. Тогда $M,N,F$ лежат на одной прямой.
По Менелаю в $ABC$:
\[
\frac{AM}{MB}\cdot\frac{BN}{NC}\cdot\frac{CF}{FA}=1.
\]
\item Подставьте $BN/NC=1/2$ и $CF/FA=5/16$:
\[
\frac{AM}{MB}\cdot\frac12\cdot\frac5{16}=1
\quad\Longrightarrow\quad
\frac{AM}{MB}=\frac{32}{5}.
\]
\end{enumerate}

\textbf{Проверка.}
Оба произведения Менелая равны единице:
\[
\frac65\cdot\frac83\cdot\frac5{16}=\frac{240}{240}=1,
\quad
\frac{32}{5}\cdot\frac12\cdot\frac5{16}=\frac{160}{160}=1.
\]

\textcolor{teal!60!black}{\textbf{Итог:}}
плоскость делит ребро $AB$ в отношении
\[\boxed{AM:MB=32:5}.\]

\textbf{Задача на закрепление.}
В пирамиде $SABC$ даны точки $P\in AS$, $Q\in CS$, $N\in CB$:
\[
AP:PS=5:3,\quad CQ:QS=4:7,\quad CN:NB=2:3.
\]
Плоскость $(PQN)$ пересекает $AB$ в точке $M$.
Найдите $AM:MB$.

\vfill
\begin{center}
{\small Лёвин Артём Александрович эксклюзивно для Ксении Клыковой}
\end{center}
\end{document}
